\begin{env}[4.3.4]
\label{0.4.3.4-fr}
Sea $\mathcal{C}$ una $\mathcal{B}$-Álgebra. Dar la estructura de
Álgebra sobre $\mathcal{C}$ equivale a dar un $\mathcal{B}$-homomorfismo
\[
  \mathcal{C}\otimes_{\mathcal{B}}\mathcal{C}\longrightarrow\mathcal{C}
\]
que satisfaga las condiciones de asociatividad y de conmutatividad
(condiciones que se verifican en cada tallo). El isomorfismo anterior
permite transportar este homomorfismo a un homomorfismo de
$\mathcal{A}$-Módulos
\[
  \Psi^*(\mathcal{C})\otimes_{\mathcal{A}}\Psi^*(\mathcal{C})
  \longrightarrow\Psi^*(\mathcal{C})
\]
que satisface las mismas condiciones. Así, $\Psi^*(\mathcal{C})$ queda
provisto de una estructura de $\mathcal{A}$-Álgebra. En particular, de
las definiciones se deduce de inmediato que la $\mathcal{A}$-Álgebra
$\Psi^*(\mathcal{B})$ es \emph{igual a} $\mathcal{A}$ (salvo isomorfismo
canónico).

Del mismo modo, si $\mathcal{M}$ es un $\mathcal{C}$-Módulo, dar esta
estructura de Módulo equivale
\oldpage[0\textsubscript{I}]{42}
a dar un $\mathcal{B}$-homomorfismo
$\mathcal{C}\otimes_{\mathcal{B}}\mathcal{M}\to\mathcal{M}$ que satisfaga
la condición de asociatividad. Por transporte se obtiene así una
estructura de $\Psi^*(\mathcal{C})$-Módulo sobre $\Psi^*(\mathcal{M})$.
\end{env}
